% GENERATED FILE. Edit canonical lessons, then run npm run generate:books.
% canonical-source-hash: fnv1a64:6c757cf45a770cbe
% canonical-lessons: RU-C313-yolka, RU-C313-girlyanda, RU-C313-ukrashat, RU-C313-ukrasheniye, RU-C313-naryazhat

\chapter{The New Year tree}
\label{ch:ru-the-new-year-tree}
\hlchaptermodality{313}

\begin{chapteropening}
\textbf{By the end of this chapter:} \emph{I can talk about a New Year tree, fairy lights, decorating, a decoration and dressing the tree.}

Complete the last lesson of chapter 313: I can talk about a New Year tree, fairy lights, decorating, a decoration and dressing the tree.
\end{chapteropening}

\section[\texorpdfstring{\ru{ёлка} (yólka)}{ёлка (yólka)}]{\ru{ёлка} (yólka) — a New Year tree; a fir tree}

\label{lesson:RU-C313-yolka}



\begin{quote}
Before the new one: say the Russian for to wish, then the Russian for a toast.
\end{quote}

\subsection*{You'll want to know: \ru{ёлка}}
\textbf{\ru{ёлка}} (\emph{yólka}) — \textquotedblleft{}a New Year tree; a fir tree\textquotedblright{}.

\begin{quote}
\ru{Перед} \ru{Новым} \ru{годом} \ru{мы} \ru{купили} \ru{ёлку}.
\end{quote}

\begin{quote}
\emph{Péred Nóvym gódom my kupíli yólku.}
\end{quote}

\begin{quote}
Before New Year we bought a tree.
\end{quote}

The Russian holiday tree, put up for New Year. \emph{\ru{Ель}} is the tree in the forest.

\begin{cousinweb}
\emph{\ru{ёлка}} (EtymDB \emph{jolka}) is built on \emph{\ru{ель}} \textquotedblleft{}spruce, fir\textquotedblright{}, inherited from Proto-Slavic \emph{jedlĭ}, with the ending \emph{-\ru{ка}}.
\end{cousinweb}

\begin{grammarlens}[title={What you've built}]
One more word for telling, explaining and describing at B1.
\end{grammarlens}

\subsection*{Your turn}
\begin{itemize}
\raggedright
  \item \emph{Say it:} \emph{yólka}
  \item \emph{Say it:} \emph{yólka}, once more
  \item \emph{Say it:} a sentence of your own with \emph{yólka}
  \item \emph{Recall:} say the Russian for to wish, then the Russian for a toast, then say \emph{yólka} again
\end{itemize}

\begin{tcolorbox}[breakable,colback=teal!4,colframe=teal!35!black,title={Before you move on}]
What does \ru{ёлка} mean? (\textquotedblleft{}A New Year tree; a fir tree\textquotedblright{}.) Say it once more.
\end{tcolorbox}
\section[\texorpdfstring{\ru{гирлянда} (girlyánda)}{гирлянда (girlyánda)}]{\ru{гирлянда} (girlyánda) — fairy lights; a garland}

\label{lesson:RU-C313-girlyanda}



\begin{quote}
Before the new one: say the Russian for a toast, then the Russian for a New Year tree.
\end{quote}

\subsection*{You'll want to know: \ru{гирлянда}}
\textbf{\ru{гирлянда}} (\emph{girlyánda}) — \textquotedblleft{}fairy lights; a garland\textquotedblright{}.

\begin{quote}
\ru{На} \ru{окне} \ru{горела} \ru{гирлянда}.
\end{quote}

\begin{quote}
\emph{Na okné goréla girlyánda.}
\end{quote}

\begin{quote}
Fairy lights were glowing in the window.
\end{quote}

Feminine. A string of lights or of paper flowers.

\begin{cousinweb}
\emph{\ru{гирлянда}} is usually traced through German or French to Italian \emph{ghirlanda} \textquotedblleft{}a garland\textquotedblright{}; EtymDB lists the Russian word with no parent.
\end{cousinweb}

\begin{grammarlens}[title={What you've built}]
One more word for telling, explaining and describing at B1.
\end{grammarlens}

\subsection*{Your turn}
\begin{itemize}
\raggedright
  \item \emph{Say it:} \emph{girlyánda}
  \item \emph{Say it:} \emph{girlyánda}, once more
  \item \emph{Say it:} a sentence of your own with \emph{girlyánda}
  \item \emph{Recall:} say the Russian for a toast, then the Russian for a New Year tree, then say \emph{girlyánda} again
\end{itemize}

\begin{tcolorbox}[breakable,colback=teal!4,colframe=teal!35!black,title={Before you move on}]
What does \ru{гирлянда} mean? (\textquotedblleft{}Fairy lights; a garland\textquotedblright{}.) Say it once more.
\end{tcolorbox}
\section[\texorpdfstring{\ru{украшать} (ukrashát)}{украшать (ukrashát)}]{\ru{украшать} (ukrashát) — to decorate}

\label{lesson:RU-C313-ukrashat}



\begin{quote}
Before the new one: say the Russian for a New Year tree, then the Russian for fairy lights.
\end{quote}

\subsection*{You'll want to know: \ru{украшать}}
\textbf{\ru{украшать}} (\emph{ukrashát}) — \textquotedblleft{}to decorate\textquotedblright{}.

\begin{quote}
\ru{Дети} \ru{украшали} \ru{ёлку} \ru{весь} \ru{вечер}.
\end{quote}

\begin{quote}
\emph{Déti ukrasháli yólku ves vécher.}
\end{quote}

\begin{quote}
The children decorated the tree all evening.
\end{quote}

Imperfective; perfective \emph{\ru{украсить}} (\emph{ukrásit}), \emph{\ru{украшу}}.

\begin{cousinweb}
\emph{\ru{украшать}} is built on \emph{\ru{у}-} and \emph{\ru{краса}} \textquotedblleft{}beauty\textquotedblright{}, an inherited Slavic word, in the shape \emph{-\ru{краш}-}: to make beautiful. EtymDB lists no parent.
\end{cousinweb}

\begin{grammarlens}[title={What you've built}]
One more word for telling, explaining and describing at B1.
\end{grammarlens}

\subsection*{Your turn}
\begin{itemize}
\raggedright
  \item \emph{Say it:} \emph{ukrashát}
  \item \emph{Say it:} \emph{ukrashát}, once more
  \item \emph{Say it:} a sentence of your own with \emph{ukrashát}
  \item \emph{Recall:} say the Russian for a New Year tree, then the Russian for fairy lights, then say \emph{ukrashát} again
\end{itemize}

\begin{tcolorbox}[breakable,colback=teal!4,colframe=teal!35!black,title={Before you move on}]
What does \ru{украшать} mean? (\textquotedblleft{}To decorate\textquotedblright{}.) Say it once more.
\end{tcolorbox}
\section[\texorpdfstring{\ru{украшение} (ukrashéniye)}{украшение (ukrashéniye)}]{\ru{украшение} (ukrashéniye) — a decoration; a piece of jewellery}

\label{lesson:RU-C313-ukrasheniye}



\begin{quote}
Before the new one: say the Russian for fairy lights, then the Russian for to decorate.
\end{quote}

\subsection*{You'll want to know: \ru{украшение}}
\textbf{\ru{украшение}} (\emph{ukrashéniye}) — \textquotedblleft{}a decoration; a piece of jewellery\textquotedblright{}.

\begin{quote}
\ru{В} \ru{коробке} \ru{лежат} \ru{ёлочные} \ru{украшения}.
\end{quote}

\begin{quote}
\emph{V koróbke lezhát yólochnyye ukrashéniya.}
\end{quote}

\begin{quote}
The tree decorations are in a box.
\end{quote}

Neuter. Also jewellery: \emph{\ru{золотые} \ru{украшения}}.

\begin{cousinweb}
\emph{\ru{украшение}} is built on \emph{\ru{украшать}} \textquotedblleft{}to decorate\textquotedblright{}, from \emph{\ru{краса}} \textquotedblleft{}beauty\textquotedblright{}, an inherited Slavic word.
\end{cousinweb}

\begin{grammarlens}[title={What you've built}]
One more word for telling, explaining and describing at B1.
\end{grammarlens}

\subsection*{Your turn}
\begin{itemize}
\raggedright
  \item \emph{Say it:} \emph{ukrashéniye}
  \item \emph{Say it:} \emph{ukrashéniye}, once more
  \item \emph{Say it:} a sentence of your own with \emph{ukrashéniye}
  \item \emph{Recall:} say the Russian for fairy lights, then the Russian for to decorate, then say \emph{ukrashéniye} again
\end{itemize}

\begin{tcolorbox}[breakable,colback=teal!4,colframe=teal!35!black,title={Before you move on}]
What does \ru{украшение} mean? (\textquotedblleft{}A decoration; a piece of jewellery\textquotedblright{}.) Say it once more.
\end{tcolorbox}
\section[\texorpdfstring{\ru{наряжать} (naryazhát)}{наряжать (naryazhát)}]{\ru{наряжать} (naryazhát) — to dress up; to trim (a tree)}

\label{lesson:RU-C313-naryazhat}



\begin{quote}
Before the new one: say the Russian for to decorate, then the Russian for a decoration.
\end{quote}

\subsection*{You'll want to know: \ru{наряжать}}
\textbf{\ru{наряжать}} (\emph{naryazhát}) — \textquotedblleft{}to dress up; to trim (a tree)\textquotedblright{}.

\begin{quote}
\ru{В} \ru{субботу} \ru{мы} \ru{нарядили} \ru{ёлку}.
\end{quote}

\begin{quote}
\emph{V subbótu my naryadíli yólku.}
\end{quote}

\begin{quote}
On Saturday we trimmed the tree.
\end{quote}

Imperfective; perfective \emph{\ru{нарядить}} (\emph{naryadít}), \emph{\ru{наряжу}}. \emph{\ru{Наряжаться}} — to dress up.

\begin{cousinweb}
\emph{\ru{наряжать}} is usually traced to \emph{\ru{на}-} and \emph{\ru{ряд}} \textquotedblleft{}a row, order\textquotedblright{}: to put in order, to array. EtymDB lists no parent.
\end{cousinweb}

\begin{grammarlens}[title={What you've built}]
One more word for telling, explaining and describing at B1.
\end{grammarlens}

\subsection*{Your turn}
\begin{itemize}
\raggedright
  \item \emph{Say it:} \emph{naryazhát}
  \item \emph{Say it:} \emph{naryazhát}, once more
  \item \emph{Say it:} a sentence of your own with \emph{naryazhát}
  \item \emph{Recall:} say the Russian for to decorate, then the Russian for a decoration, then say \emph{naryazhát} again
\end{itemize}

\begin{tcolorbox}[breakable,colback=teal!4,colframe=teal!35!black,title={Before you move on}]
What does \ru{наряжать} mean? (\textquotedblleft{}To dress up; to trim (a tree)\textquotedblright{}.) Say it once more.
\end{tcolorbox}
